\documentclass[11pt]{article}
% Standalone preamble for web-only supplementary notes.
\usepackage[T1]{fontenc}
\usepackage{amsmath,amssymb,amsthm}
\usepackage{hyperref}
\usepackage{cleveref}
\DeclareMathOperator{\Map}{Map}
\DeclareMathOperator{\Fun}{Fun}
\newcommand{\iso}{\xrightarrow{\sim}}
\newtheorem{lemma}{Supplementary Lemma}
\newtheorem{corollary}[lemma]{Supplementary Corollary}
\crefname{lemma}{Supplementary Lemma}{Supplementary Lemmas}
\Crefname{lemma}{Supplementary Lemma}{Supplementary Lemmas}
\crefname{corollary}{Supplementary Corollary}{Supplementary Corollaries}
\Crefname{corollary}{Supplementary Corollary}{Supplementary Corollaries}

\begin{document}
\section*{Triangles over a base}

The definition of \href{exponentiable-functors.html}{functors over a base}
remembers the identification making the triangle commute. This note explains
how that identification is retained when passing to the relative mapping anima.
Assume the Chapter 1 structures, including mapping animae, functor categories
and pullbacks. All categories and functors below are absolute.

Fix $f\colon C\to S$ and $g\colon D\to S$. Write a functor over $S$ as
$u=(h,\theta)$, where $h\colon C\to D$ and $\theta\colon gh=f$.
We orient the triangle towards $f$, as in the code; inverting $\theta$
gives the convention in the book. For $v=(k,\psi)$, a comparison $u\simeq_S v$
consists of $\alpha\colon h=k$ and a specified identification
\[
 \psi\cdot(g\alpha)=\theta.
\]
Here $g\alpha$ denotes postcomposition of the identification by $g$, and
$\beta\cdot\alpha$ means first $\alpha$, then $\beta$.

\begin{lemma}
\label[lemma]{lem:Supplementary_Relative_Names}
A triangle $u$ has a name $\ulcorner u\urcorner$ in $\Map_S(C,D)$.
A point $x$ of this anima determines a triangle $\operatorname{decode}(x)$.
There are comparisons
\[
 \ulcorner\operatorname{decode}(x)\urcorner=x,
 \qquad \operatorname{decode}(\ulcorner u\urcorner)\simeq_S u.
\]
\end{lemma}
\begin{proof}[Proof idea]
The category $\Fun_S(C,D)$ is the pullback of
$g_*\colon\Fun(C,D)\to\Fun(C,S)$ along the point naming $f$.
Name $h$ in the first factor and use the identity of the terminal category
in the second. Naming $\theta$, followed by the postcomposition and unit
comparisons, supplies the matching of this cone. The map converting
$\theta$ into that matching is an equivalence of identification animae.

Conversely, decode the first leg of a point of the pullback, normalize its
terminal leg, and lift its matching through this equivalence. The pullback
computation and reconstruction comparisons give the two displayed formulas.
For the second formula, the matching comparison also recovers the specified
triangle. Passing to the core gives the statement for $\Map_S(C,D)$.
\end{proof}

\begin{lemma}
\label[lemma]{lem:Supplementary_Relative_Comparisons}
A comparison $u\simeq_S v$ gives an identification
$\ulcorner u\urcorner=\ulcorner v\urcorner$. Conversely, an identification
of these names gives a comparison $u\simeq_S v$.
\end{lemma}
\begin{proof}[Proof idea]
Name the underlying identification $\alpha$. Its equation over $S$ becomes
the matching equation for a comparison of the two named pullback cones.
Lift that whole cone comparison and then name the result in the core.

In the reverse direction, extract a cone comparison from the identification
of names. Its terminal component is identified with the identity. Lift the
other component through the naming equivalence. Reflecting the matching
equation through the equivalence used in \Cref{lem:Supplementary_Relative_Names}
gives exactly $\psi\cdot(g\alpha)=\theta$.
These constructions provide comparisons in both directions; the linked
results here do not identify the two operations on comparison data as
mutual inverses with their specified higher witnesses.
\end{proof}

\begin{lemma}
\label[lemma]{lem:Supplementary_Relative_Associator}
Compose functors over $S$ by pasting their triangles, with the associator
inserted at the composite. The ordinary associator and the left and right
unit identifications then give comparisons over $S$.
\end{lemma}
\begin{proof}[Proof idea]
For $u=(h,\theta)\colon f\to g$ and $v=(k,\psi)\colon g\to l$, the
composite triangle follows the route
\[
 l(kh)\simeq(lk)h\xrightarrow{\psi h}gh\xrightarrow{\theta}f.
\]
For three functors, the two pasted routes are compared by the
projection-square pentagon calculation. The unit cases use the corresponding
triangle and unit calculations. These prove compatibility of the displayed
associator and unit identifications with the triangles over $S$.
A pentagon among these relative associators is a further assertion, not part
of this lemma.
\end{proof}
\end{document}
